\subsubsection{Completing a phase}
\label{ha:postprocess}

We complete the induction in Lemma~\ref{ha:construction} by handling
the border balls.  At the end of processing the bags of
$\mathcal F_\ell$, let $\operatorname{Bor}_t$ and
$\operatorname{Int}_t$ be the unprocessed border and interior balls
left by its loop.  The loop stops with
$|\operatorname{Int}_t|<|\operatorname{Bag}_t|/2$.
It is enough to give each $t\in T_\ell$ an additional
\[
 \frac{C_4}{6C_3}|\operatorname{Bor}_t|\delta_\ell
\]
tokens.  Together with the credit already obtained during its loop,
this gives at least
\begin{equation}\label{ha:phase-target}
 \frac{C_4}{6C_3}
 (|\operatorname{Bag}_t|-|\operatorname{Int}_t|)\delta_\ell
 \ \ge\ \frac{C_4}{12C_3}|\operatorname{Bag}_t|\delta_\ell.
\end{equation}

Write $Z$ for the collection just before this last step,
$p(B)=\operatorname{tok}_\ell(B)$, and let $q(B)$ be the token
value of the same unchanged record at the beginning of phase
$\ell+1$.  Thus for an original ball inserted in the current phase
both $p(B)$ and $q(B)$ equal $\delta_\ell-C_4\delta_{\ell+1}$.
For all other nonavoiding records the threshold in the positive
part decreases from $C_4\delta_\ell$ to $C_4\delta_{\ell+1}$;
for avoiding records the same formula is divided by $2(C_4+2)$.
In particular $q(B)\ge p(B)$.
Set
\[
 b=\sum_{t\in T_\ell}|\operatorname{Bor}_t|,
 \qquad e=\sum_{B\in Z}(q(B)-p(B)).
\]
If $b=0$ there is no border credit to supply, and the phase is
complete.  If $e\ge (C_4/(6C_3))b\delta_\ell$, retain $Z$ and
distribute these extra tokens among the current nodes in proportion
to $|\operatorname{Bor}_t|$.  Every node obtains the required
additional amount, and all earlier credits remain available.

It remains to handle
\begin{equation}\label{ha:small-excess}
 e<\frac{C_4}{6C_3}b\delta_\ell.
\end{equation}
Leave current-phase original balls unchanged.  Replace every other
record $B$, with center $x$ and numerical radii $r_1<r_2$, by the
concentric restriction
\[
 S(B)=B\bigl(x,r_1+(C_3+1)\delta_\ell\ \Vert\
                    r_2-(C_3+1)\delta_\ell\bigr),
\]
discarding it if the displayed interval has nonpositive width.
Use the same labels on each surviving restriction, and assign a
discarded record token value zero.  Finally insert all balls of
$\bigcup_{t\in T_\ell}\operatorname{Bor}_t$ with their original
labels.

Here is a direct verification of separation, including the role
of conflict domains.  For a border ball $A=B(X_u,\delta_\ell)$,
consider an old record $B$ whose conflict domain meets $C(A)$.
Lemma~\ref{ha:border-monotone} says that $A$ is still border, so
\[
 \|X_u-x\|_1\le r_1+C_3\delta_\ell
 \quad\hbox{or}\quad
 \|X_u-x\|_1\ge r_2-C_3\delta_\ell.
\]
In the first case every point of $A$ has distance strictly less
than the new inner radius; in the second it has distance strictly
greater than the new outer radius.  Thus $A$ does not meet $S(B)$.
If the conflict domains are disjoint no geometric separation is needed.
The original balls in the current family are pairwise disjoint,
so no conflict arises either with an unchanged current-phase ball
or with another newly inserted border ball.  The other validity
conditions follow from Lemmas~\ref{ha:input-valid} and
\ref{ha:heredity}.  Hence the new collection is valid.

\begin{lemma}[The cost of shrinking]\label{ha:shrink-cost}
Under \eqref{ha:small-excess}, the total token deficit created by
the restrictions satisfies
\[
 \mathsf D:=\sum_{B\in Z}\bigl(p(B)-q(S(B))\bigr)
 \le\frac12 b\delta_\ell.
\]
Every summand defining $\mathsf D$ is nonnegative.
\end{lemma}
\begin{proof}
For unchanged original balls both the deficit and $q(B)-p(B)$
are zero.  For any other record the decrease of width is
$2(C_3+1)\delta_\ell$, unless the record disappears entirely.
This exceeds the decrease of the threshold, so the deficit is
nonnegative.  We prove, record by record,
\begin{equation}\label{ha:shrink-comparison}
 q(B)-p(B)\ge\frac{C_4}{3C_3}
                   \bigl(p(B)-q(S(B))\bigr).
\end{equation}
If $p(B)=0$, its deficit is zero and the inequality is immediate.
Otherwise, first take a nonavoiding record of width $w$.
Since $w>C_4\delta_\ell$ and
$\delta_{\ell+1}\le\delta_\ell/3$,
\[
 q(B)-p(B)=C_4(\delta_\ell-\delta_{\ell+1})
             \ge\frac23 C_4\delta_\ell.
\]
Using the positive-part formula, whether or not the restriction
survives,
\[
 p(B)-q(S(B))
 \le2(C_3+1)\delta_\ell-C_4(\delta_\ell-\delta_{\ell+1})
 \le2C_3\delta_\ell,
\]
where the last step uses $C_4\ge3$.
These two bounds imply \eqref{ha:shrink-comparison}.
For avoiding records both sides of these estimates have the common
factor $1/(2(C_4+2))$, giving exactly the same conclusion.
Summing \eqref{ha:shrink-comparison} and then using
\eqref{ha:small-excess} gives
$\mathsf D\le(3C_3/C_4)e\le b\delta_\ell/2$.
\end{proof}

Each newly inserted border ball has
$\delta_\ell-C_4\delta_{\ell+1}$ tokens.  Reserve at most
$\delta_\ell/2$ of these to cover the deficits of the old records.
Lemma~\ref{ha:shrink-cost} guarantees enough in the resulting pool:
if $\mathsf D>0$, distribute the pool in proportion to the individual
deficits, stopping once each is paid; if $\mathsf D=0$, no transfer is
needed.  Send the tokens of $S(B)$ together with this payment back
to the old record $B$, and redistribute its former credits by
induction.  Each border ball still has at least
\[
 \delta_\ell-C_4\delta_{\ell+1}-\frac12\delta_\ell
 \ge\frac13\delta_\ell
 \ge\frac{C_4}{6C_3}\delta_\ell
\]
tokens for its own node.  The first inequality uses
$C_4\lambda\le1/6$; the second follows from the stated lower bound
on $C_3$.  This supplies precisely the missing border credit.

Both phase-transition cases preserve all previous credits and
give \eqref{ha:phase-target} to every node of the current family.
Starting with $Z=\varnothing$ and applying the interior construction
and the phase transition for $\ell=1,\ldots,m$ therefore yields
\[
 \frac{C_4}{12C_3}
 \sum_{i=1}^m\delta_i\sum_{t\in T_i}|\operatorname{Bag}_t|
 \le\sum_{B\in Z}\operatorname{tok}_{m+1}(B)
 \le\sum_{B\in Z}\operatorname{width}(B).
\]
This proves Lemma~\ref{ha:insertable-arrangement}, completing the
proof of Proposition~\ref{prop:hierarchy-assigned} by the reduction
above.  If a countable geometric sequence is used, apply the finite
statement to every initial segment and take the increasing limit
of its nonnegative right-hand mass; no limiting arrangement is
required.
